\section{现状与痛点}

\subsection{基层经营管理场景中的重复性资料处理负担}

在片区综合岗位的日常工作中，经营资料的收集、甄别、重命名与归档等基础性事务占据了同事们的部分工作精力。此类事务单次操作复杂度并不高，但频次高、覆盖面广、格式差异大，因而形成了持续存在的隐性工作负担。以月度综合财务检查、节假日值班安排、人员考勤汇总、技能鉴定统计以及库存盘点资料为例，上报文件虽来源不同、命名方式不一，但通常都遵循较为稳定的业务语义结构，即文件名中往往包含站点名称、业务主题和时间属性等关键字段。

例如，综合财务检查资料的文件命名，通常由“站点名称”“财务合规检查问题汇总表”等核心字段，与少量个性化补充表述共同构成；节假日值班资料通常体现为“站点名称+节假日名称+值班表”的结构；考勤等资料也普遍具有“站点名称+业务主题”的共同规律。再特殊一些的库存盘点资料，虽然往往以目录而非单一文件的形式出现，但其内部通常仍包含核对表、盘点表、临期商品台账、异常库存明细等子项，这些文件同样没有脱离“关键字+站点名称”的基本命名逻辑。即使个别文件命名不够规范，相关语义信息仍可从文件标题、正文内容甚至附属元数据中被识别出来。这表明，基层经营资料管理虽然表面上呈现为非标准化问题，但其底层实际上具有较强的规则性、可抽取性与可自动化潜力。

\subsection{经营分析资料处理中的认知负荷问题}

除资料归档外，片区综合岗位还需要定期处理大批量经营分析数据。例如，在月度或季度非油业务分析过程中，工作人员往往需要同时面对商品大类库存明细、收入毛利明细、劳动竞赛数据以及其他经营报表。这些文件来源多元、体量较大，单个文件规模甚至可接近 \SI{50}{\mega\byte}。在连续查阅、筛选和比对大量表格之后，工作人员不仅需要提炼关键指标，还需进一步形成综合性业务简报，其工作过程具有明显的高重复性与高注意力消耗特征。

从认知机理看，这种主观体验并非个体化感受，而具有明确的研究依据。相关研究指出，长时间持续执行单调且高强度的信息处理任务，会引发心理疲劳，并表现为反应时间延长、正确率下降以及对关键信息的注意资源分配减少\cite{guo2016mental,oken2006vigilance}。这与基层经营分析中的实际体验高度一致，即在长时间面对结构相近、指标密集的数据报表时，工作人员容易出现注意稳定性下降、信息辨识效率降低以及短时理解迟滞等现象。由此可见，当前经营分析工作的主要瓶颈不仅来自业务量本身，更来自“重复识别+持续比较+人工归纳”所叠加产生的认知负荷。

\subsection{规则稳定性为智能化介入提供了现实基础}

值得注意的是，此类工作虽然耗时耗力，但并非不可描述、不可建模。无论是资料命名规则，还是月报、季报中反复出现的字段结构、指标口径和分析模板，均具有较强的稳定性。大多数经营分析任务在不同周期之间的差异，主要体现在数据取值、文件命名和局部异常点上，而非业务逻辑的根本变化。因此，这类任务天然适合通过程序化工具或 AI 智能体进行重构。

从流程视角看，若能够将“识别文件类型”“提取关键字段”“按统一规则重命名”“汇总核心指标”“生成分析初稿”等步骤转化为可执行的规则链条，则原本依赖人工经验与反复操作的事务性工作，就可以被迁移为“人制定规则、智能体执行流程、人工复核结果”的新型协同模式。这种模式并不取代业务人员的专业判断，而是将其从低附加值、强重复性的劳动中释放出来，使其更聚焦于异常识别、经营诊断和管理决策支持等更具价值的工作环节。

\subsection{自然语言交互拓展了办公自动化边界}

除文件归集与报表分析外，基层办公场景还广泛存在另一类典型问题，即工具功能可用但使用门槛较高。以 Excel 为例，工作人员在处理数据时常常会出现函数语法遗忘、函数选型不确定或复杂嵌套公式难以构造等情况。传统解决方式通常依赖经验查找、网络搜索或反复试错，效率有限，且容易打断原有工作节奏。

随着大语言模型技术的发展，自然语言到结构化表达的转换能力已在办公场景中展现出明显潜力。Zhao 等提出的 NL2Formula 研究表明，大语言模型及相关方法能够根据自然语言查询生成可执行的电子表格公式，从而为复杂函数编写提供新的技术路径\cite{zhao-etal-2024-nl2formula}。这意味着，基层员工可以逐步从“记忆函数语法”的工作方式，转向“描述业务意图，由模型辅助生成公式”的交互方式。对经营管理岗位而言，这不仅提升了办公软件的使用效率，也进一步说明 AI 智能体在资料归集、数据整理与辅助分析等场景中具有现实可行性和推广价值。
